\section{Implementation in the BX3 Framework}\label{sec:implementation}

The preceding sections specify the Bailout Protocol as a formal state machine
$\mathcal{B} = (Q, q_0, \Sigma, \rightarrow, Q_F)$ and as a compulsory escalation
mechanism.  This section describes how that specification is realized within the
\bxthree{} three-layer architecture.  We show precisely how each layer
participates in protocol execution, how the \bounds{} Engine enforces the
trigger condition, how the \fact{} Layer provides the enforcement substrate and
maintains the append-only Forensic Ledger, and how the state machine is embedded
as a deterministic \fact{} Layer extension that makes the upstream accountability
guarantee structurally operative at runtime.

% --------------------------------------------------------------
\subsection{Purpose Layer: Human Accountability Anchor}\label{sec:purpose-layer}
% --------------------------------------------------------------

The \purpose{} layer carries intent, judgment, and final authority.  Within the
Bailout Protocol, its role is singular and architecturally irreplaceable: it is
the \emph{exclusive terminus} for all exception propagation, and it is the only
layer capable of issuing a binding resolution directive that closes a protocol
run.

Each \bxthree{} node $n$ is provisioned at creation time with a reference to a
named human principal $h(n)$.  This reference is stored in the node's
\purpose{} layer worksheet and is immutable for the node's operational
lifetime.  No machine actor---including the node's own \bounds{} Engine, any
parent node, or any recursively spawned subordinate---may modify $h(n)$ without
explicit human authorization recorded as a ledger entry.  The immutability of
$h(n)$ is the structural fact that terminates the escalation chain and makes
the upstream accountability guarantee unconditional.

The \purpose{} layer's participation in the Bailout Protocol proceeds through
four discrete functions:

\begin{itemize}\itemsep=0pt
\item[(\textit{i})] \textbf{Provisioning.}  At node initialization, the
\purpose{} layer registers $h(n)$ with the \fact{} Layer's Pathway Registry and
records the anchor assignment in the authorization ledger.  This entry
constitutes the human principal's explicit acceptance of accountability for node
$n$'s operational scope.  The anchor is non-negotiable and non-migratory: for
all nodes $n_i$ in any spawn chain $C$ descended from $n$, $h(n_i) = h(n)$.

\item[(\textit{ii})] \textbf{Exception receipt.}  When an exception state
propagates to the \purpose{} layer via the Escalate algorithm, the \purpose{}
layer receives the full diagnostic context accumulated along the propagation
chain: the originating node, the trigger condition $\text{TRIGGER}(n,x)$, the
operating range definition $R(n)$, the time since first detection, and any
prior resolution history for this exception class.

\item[(\textit{iii})] \textbf{Directive issuance.}  The \purpose{} layer
evaluates the exception and issues one of three binding directives:
\texttt{ACK} (acknowledge and resume normal operation), \texttt{RESOLVE}
(execute a specified binding resolution and resume), or \texttt{REJECT}
(disallow the proposed action and enter a quiescent error state).  No machine
actor may issue, simulate, or substitute any of these directives.  Upon
receipt of the directive, the state machine transitions to
\texttt{HUMAN\_ACKNOWLEDGED} and then to \texttt{RESOLVED}; the directive is
committed to the authorization ledger as a \purpose-layer-authorized entry.

\item[(\textit{iv})] \textbf{Consequential authority.}  The \purpose{} layer is
the only layer that carries the authority to commit consequential decisions to
the authorization ledger.  This exclusivity makes it the \emph{only} permissible
terminus for the accountability chain: any attempt by a machine actor to close a
protocol run without a ledger entry bearing a \purpose-layer-authorized
directive fails the \fact{} Layer's pre-commit hook.
\end{itemize}

The human anchor $h(n)$ is not a design option.  It is a structural necessity.
The functional separation between \purpose{}, \bounds{}, and \fact{} layers
establishes that the \purpose{} layer is the only layer carrying both intent and
final accountability.  The \bounds{} Layer detects exception conditions but
carries no judgment authority; the \fact{} Layer enforces transition relations
and records events but carries no intent authority.  When an exception exceeds the
\bounds{} Layer's provisioned scope, no machine actor exists that has the
authority to adjudicate it.  The decision must return to the \purpose{} layer,
which carries the intent and accountability that any such adjudication requires.
The Bailout Protocol is the mechanism by which this structural necessity is
enforced at runtime.

% --------------------------------------------------------------
\subsection{Bounds Engine: Triggering Bail-Out}\label{sec:bounds-trigger}
% --------------------------------------------------------------

The \bounds{} layer carries bounded reasoning and constrained execution within
the provisioned operating range $R(n)$.  Its participation in the Bailout
Protocol is confined to a single, non-discretionary function: detecting
conditions that exceed $R(n)$ and triggering the state machine's compulsory
transition to \texttt{BOUNDED}.

The \bounds{} Engine evaluates the trigger condition $\text{TRIGGER}(n,x)$
against every input element $x$ in the node's observable input space.  This
evaluation is a read-only query against the operating range definition $R(n)$,
not a judgment call made by a machine actor.  When all three sub-conditions of
TC hold simultaneously---$x \notin R(n)$, the condition is not previously
resolved, and $\text{confidence}(n,x) < \tau$---the trigger fires and the state
machine transitions from \texttt{IDLE} to \texttt{BOUNDED} via T1.

Critically, the \bounds{} Engine exercises no discretion over this transition:

\begin{enumerate}\itemsep=0pt
\item The evaluation of $\text{TRIGGER}(n,x)$ is driven directly by the Bailout
Monitor's pre-commit hook, which intercepts the query and commits the result to
the state machine before any agent-level code can act on the input.  There is
no intervening decision step at which a machine actor could suppress or
reclassify the exception.

\item The \bounds{} Engine may not delay, reclassify, or absorb an exception
once $\text{TRIGGER}(n,x)$ evaluates to true.  The state machine's deterministic
transition relation is the only permissible response to a satisfied trigger
condition.

\item The confidence threshold $\tau$ is provisioned at node initialization and
is not adjustable at runtime by any machine actor.  This prevents a machine
actor from inflating its own confidence to suppress a warranted escalation ---
a specific failure mode the no-discretion property addresses.
\end{enumerate}

The secondary trigger---the explicit \texttt{bailout\_signal}---handles the case
where the \bounds{} Engine detects a logical contradiction or operational
impossibility that it recognizes as unresolvable even though both contributing
inputs fall individually within $R(n)$.  The \texttt{bailout\_signal} bypasses
TC entirely and immediately initiates T1, ensuring that the trigger is not
limited to conditions expressible as simple boundary violations.

Upon entering \texttt{BOUNDED}, the \bounds{} Engine is granted $T_{\text{bounds}}$
--- a provisioned time window --- to attempt local resolution.  If the
\bounds{} Engine produces a resolution $r$ that the \fact{} Layer approves before
$T_{\text{bounds}}$ expires, the state machine returns to \texttt{IDLE} via T2 and
normal execution resumes.  If all resolution paths are exhausted or
$T_{\text{bounds}}$ expires, the state machine advances to \texttt{BAIL\_PENDING}
via T3 and immediately, atomically, to \texttt{ESCALATING} via T4 --- with no
dwell time and no machine-actor intervention permitted between
\texttt{BAIL\_PENDING} and \texttt{ESCALATING}.

This atomic transition is the architectural guarantee that the bailout trigger
is compulsory.  The \bounds{} Engine cannot suppress the escalation at the moment
of maximum temptation---when $T_{\text{bounds}}$ expires and the exception must
propagate.  The \fact{} Layer's pre-commit hook enforces the transition relation
directly, bypassing any agent-level code that might otherwise interpose.

% --------------------------------------------------------------
\subsection{Fact Layer: Enforcement Substrate}\label{sec:fact-layer}
% --------------------------------------------------------------

The \fact{} layer carries deterministic enforcement and forensic auditability.
Within the Bailout Protocol, it serves two distinct but cooperating functions:
enforcing the state machine's transition relation, and maintaining the
append-only Forensic Ledger as an immutable record of every protocol event.

The \fact{} Layer is the only mechanism in the \bxthree{} architecture capable
of authorizing a state transition or committing an entry to the authorization
ledger.  This exclusivity is what makes its enforcement structurally reliable.
No machine actor---including the \bounds{} Engine, any spawned subordinate
agent, or any external system operating outside the \fact{} Layer's
jurisdiction---can authorize a transition that violates the state machine's
transition relation or commit a ledger entry that misrepresents the protocol's
actual execution.

The \fact{} Layer's enforcement of the Bailout Protocol proceeds through three
cooperating mechanisms:

\begin{itemize}\itemsep=0pt
\item[(\textit{i})] \textbf{Pre-commit hook.}  Before any state transition is
committed, the \fact{} Layer's Bailout Monitor evaluates whether the proposed
transition satisfies the current state's enabled transition set.  If no
transition is defined for the $(q,\sigma)$ pair, the transition is rejected and a
protocol violation is recorded in the Forensic Ledger.  This hook runs
synchronously and atomically: no intermediate or partially-executing state is
observable externally.

\item[(\textit{ii})] \textbf{Authorization ledger gate.}  The \fact{} Layer is
the sole write authority for the authorization ledger.  Any resolution that
purports to close a protocol run without a corresponding \purpose-layer-authorized
directive fails the gate and is rejected.  This prevents the \bounds{} Engine
from issuing a retroactive resolution after timeout to suppress a warranted
escalation.

\item[(\textit{iii})] \textbf{Pathway Registry.}  The \fact{} Layer maintains the
Pathway Health Registry, tracking all communication pathways to the human anchor
$h(n)$.  When a pathway is degraded or unavailable, the registry enables the
Escalate algorithm to identify and route through alternative functional
pathways.  Historical health data is append-only; current-pathway state is a
mutable field updated by the Bailout Monitor.
\end{itemize}

% --------------------------------------------------------------
\subsection{Forensic Ledger Recording}\label{sec:forensic-ledger}
% --------------------------------------------------------------

The Forensic Ledger is the immutable evidentiary record produced by the
\fact{} Layer's enforcement of the Bailout Protocol.  Every protocol event---
trigger evaluations, state transitions, pathway selections, acknowledgments,
directives, and timeout conditions---is committed to the ledger as an
append-only entry.  Each entry is timestamped, attributed to the originating
node, and linked to its predecessor by a SHA-256 hash chain that makes
retrospective tampering computationally detectable.

The Forensic Ledger records each protocol event across three discrete
accountability planes:

\begin{enumerate}\itemsep=0pt
\item \textbf{Purpose Plane.}  Records the authorization context at the time
of bailout---the originating node, its human anchor $h(n)$, the active mandate,
and the constraint boundary or trigger condition that was exceeded.  This plane
establishes the provenance of the exception in the delegation chain.

\item \textbf{Bounds Plane.}  Records the diagnostic context accumulated
during propagation---the trigger condition, the proposed resolution attempt, the
\bounds{} Engine's confidence level at time of firing, and the complete
propagation chain showing every node the exception traversed before reaching
$h(n)$.

\item \textbf{Fact Plane.}  Records the physical or digital system state at the
time of bailout---actuator positions, sensor readings, system telemetry, and any
physical consequences of actions taken or blocked prior to bailout resolution.
This plane provides the evidentiary baseline against which post-resolution
consequences can be measured.
\end{enumerate}

The Forensic Ledger serves three distinct accountability functions:

\begin{enumerate}\itemsep=0pt
\item \textbf{Attribution.}  Every escalation is attributable to a specific
node, trigger condition, and human anchor.  The chain of diagnostic context
accumulated during propagation is preserved in full, enabling post-hoc
reconstruction of the escalation's origin and propagation path.

\item \textbf{Non-repudiation.}  The human principal's directive ---
\texttt{ACK}, \texttt{RESOLVE}, or \texttt{REJECT} --- is immutably recorded as a
\purpose-layer-authorized ledger entry.  The principal cannot credibly deny
having issued the directive, and the system cannot fabricate a directive that
was not issued.

\item \textbf{Verifiability.}  The SHA-256 hash chain enables any authorized
auditor to verify that the ledger has not been altered since the protocol run.
Each entry's hash covers the entry's content and the hash of the preceding
entry, making insertion, deletion, or modification of historical entries
computationally impractical.
\end{enumerate}

The Forensic Ledger is not a conventional audit log.  It is a structurally
append-only data structure whose immutability is enforced by the \fact{} Layer's
exclusive write authority and by the hash-chain integrity mechanism.  A machine
actor cannot delete, modify, or suppress a ledger entry without detection.

% --------------------------------------------------------------
\subsection{State Machine as Fact Layer Extension}\label{sec:state-machine-fact-layer}
% --------------------------------------------------------------

The Bailout Protocol state machine $\mathcal{B}$ is embedded in the \fact{}
layer of every \bxthree{} node as a \fact{} Layer extension --- a deterministic
enforcement automaton whose transition relation is enforced by the \fact{} Layer's
pre-commit hook, whose state is stored in the \fact{} Layer's worksheet, and whose
state transitions generate entries in the Forensic Ledger.

This embedding is not incidental.  The state machine's enforcement requirements
are precisely matched to the \fact{} Layer's architectural properties:

\begin{itemize}\itemsep=0pt
\item[(\textit{i})] \textbf{Deterministic transition enforcement.}  For every
$(q,\sigma) \in Q \times \Sigma$, at most one transition is defined.  This
determinism is delivered by the Bailout Monitor's priority function $\pi$, which
totally orders simultaneous trigger conditions before committing any state
change.  The \fact{} Layer provides the scheduling guarantee and atomic
pre-commit hook required to make this ordering consistent across all nodes.

\item[(\textit{ii})] \textbf{Atomic state transitions.}  No intermediate or
partially-executing state is observable externally.  The \fact{} Layer's
pre-commit hook evaluates and commits the transition in a single synchronized
operation, preventing any interleaving machine actor from observing or modifying
the state mid-transition.

\item[(\textit{iii})] \textbf{State observability.}  The state machine's
current state must be observable by the parent node and by the Bailout Monitor
without granting write access.  The \fact{} Layer exposes a state observation
interface to the parent node, enabling parent-level monitoring of child
escalation status independently of the child node's internal representation.

\item[(\textit{iv})] \textbf{Ledger co-commitment.}  Every state transition
generates a mandatory Forensic Ledger entry as a co-commitment of the transition
itself.  No state change in $\mathcal{B}$ is committed without a corresponding
ledger entry recording the transition event, the triggering condition, and the
diagnostic context.
\end{itemize}

The state machine is therefore not merely \emph{located} in the \fact{} Layer---
it is defined by the \fact{} Layer's enforcement capabilities.  The six-state
machine, the nine-transition relation (T1--T9), the timeout parameters, and the
escalation algorithm are all realized as \fact{} Layer primitives.  The
\fact{} layer does not merely host the state machine; it \emph{is} the state
machine's enforcement substrate.

The three state invariants maintained by the \fact{} Layer's enforcement
architecture are:

\begin{align}
\text{INV-1:}&\quad \forall n \in \text{nodes},\; \text{state}(n) \in
Q \setminus \texttt{ESCALATING} \iff \neg\text{active\_bailout}(n).
\end{align}

INV-1 is enforced by the \fact{} Layer's pre-commit hook: a node may enter a
non-escalating state only when the Bailout Monitor confirms that no active
propagation is in flight.  The state and the active-bailout flag are updated
atomically, preserving the equivalence.

\begin{align}
\text{INV-2:}&\quad \text{state}(n) = \texttt{BAIL\_PENDING} \implies
\text{active\_bailout}(n) \land
\text{transition}(\texttt{BAIL\_PENDING}, e_{\text{timeout}}) = \texttt{ESCALATING}.
\end{align}

INV-2 follows from the \fact{} Layer's enforcement of the transition relation.
Any node in \texttt{BAIL\_PENDING} has an active bailout, and the only transition
defined for this state with an enabled event is to \texttt{ESCALATING}.  The
Bailout Monitor enforces this transition compulsorily upon receipt of
$e_{\text{timeout}}$; no machine actor may interpose.

\begin{align}
\text{INV-3:}&\quad \text{state}(n) = \texttt{RESOLVED} \implies
\exists\; d \in \{\texttt{ACK},\texttt{RESOLVE},\texttt{REJECT}\}:\;
d \in \text{ledger} \land \text{directive\_source}(d) = h(\text{originating\_node}(n)).
\end{align}

INV-3 enforces the non-negotiable requirement that every protocol run that
reaches \texttt{RESOLVED} has a corresponding \purpose-layer-authorized
directive from the human anchor $h(n)$ of the originating node.  The directive
must exist in the ledger and must be sourced from $h(n)$, not from any machine
actor.  This invariant is checked by the \fact{} Layer's ledger gate at the
time of the \texttt{RESOLVED} transition and is the architectural expression of
the no-machine-actor-terminus property.

Together, the state invariants and the \fact{} Layer's enforcement architecture
ensure that the Bailout Protocol state machine operates as a deterministic,
atomic, auditable extension of the \fact{} Layer---and that every transition it
executes is immutably recorded in the Forensic Ledger as evidence that the
upstream accountability guarantee has been upheld for the current protocol run.

% --------------------------------------------------------------
\subsection{Relationship: Bailout Protocol as Runtime Expression of the
Upstream Accountability Guarantee}\label{sec:runtime-expression}
% --------------------------------------------------------------

The upstream accountability guarantee of the \bxthree{} Framework is the
commitment that \bxthree{} systems \emph{fail upward into human consciousness,
never downward into autonomous chaos}.  The Bailout Protocol is the runtime
mechanism that makes this guarantee structurally operative---not for the
exception conditions anticipated and resolved by the four preceding pillars, but
for the complementary set of all unanticipated exceptions that exceed
machine-actor provisioned scope.

The relationship between the protocol and the guarantee is one of mutual
structural definition.  The upstream accountability guarantee requires the Bailout
Protocol to be operative: without it, the guarantee would be unenforceable for
the exception conditions that the four preceding pillars do not cover.  The
guarantee is a claim about the behavior of the system; the protocol is the
mechanism that makes the claim true.  Conversely, the Bailout Protocol requires
the structural properties that only the \bxthree{} three-layer architecture
provides: the \purpose{} layer must be non-delegable as an exception terminus,
the $h(n)$ anchor must be immutable at runtime, and the \fact{} layer must have
exclusive write authority over the Forensic Ledger.  These are not optional
features that improve the protocol's performance---they are structural
requirements without which the protocol's guarantees cannot be realized.

The structural completeness of this three-layer implementation is what makes the
upstream accountability guarantee unconditional.  The guarantee does not say that
human oversight is available if machine actors choose to escalate; it says that
human oversight is \emph{compulsory} for every unresolvable exception, regardless
of machine-actor preferences, operational urgency, or the availability of
alternative resolution paths.

The five pillars of \bxthree{} collectively establish the conditions under
which the upstream accountability guarantee applies.  Loop Isolation, Recursive
Spawning, Spatial Firewall, and Sandbox Gate address exception conditions
anticipated at design time and resolvable within the machine-only execution
graph.  The Bailout Protocol addresses the complementary remainder---conditions
not anticipated, that exceed the operating range of any machine actor, and that
cannot be resolved within the machine-only execution graph.  The completeness of
this coverage is what makes the upstream accountability guarantee unconditional:
it is not that \emph{most} exceptions reach $h(n)$, but that \emph{every}
exception reaches $h(n)$, because the trigger condition fires on any condition
outside $R(n)$ and the escalation path has no machine-actor terminus by
architectural constraint.

This unconditionality is what distinguishes the \bxthree{} upstream
accountability guarantee from all prior escalation architectures.  Tiered
Agentic Oversight (TAO) models permit human oversight to function as a
high-tier option at the top of a machine-tier hierarchy, where each intermediate
tier may absorb the exception before passing it up.  The human anchor in TAO is
present but not architecturally required for every exception; absorption at an
intermediate tier is permitted and invisible to external auditors.  The Bailout
Protocol forbids this absorption by architectural constraint: the human anchor
is not one option among many---it is the \emph{only} permissible terminus for
every unresolved exception.  The parent-chain bypass skips intermediate nodes
not solely for efficiency but because their presence in the terminus would break
the accountability chain by introducing machine actors as potential terminal
nodes.

The \fact{} Layer's state machine extension ensures that the upstream
accountability guarantee is not merely a design claim but a runtime observable
property.  Every state transition is recorded; every protocol run produces a
Forensic Ledger entry bearing the human anchor's directive; every invariant
check is enforced by the \fact{} Layer's pre-commit hook.  The guarantee is not
reconstructed post-hoc from logs --- it is established by the architecture at
the moment of each protocol run.